% !TEX root = ../tecnologia-en-marcha.tex
% !TEX encoding = UTF-8 Unicode
% Texto: primer borrador del artículo, con acentos restaurados y citas enlazadas.

\section{Resultados}

El marco fue ejecutado sobre dos configuraciones piloto. El
Cuadro~\ref{cuadro:resultados} sintetiza los resultados generales del
expediente.

\begin{table}[htbp]
  \caption{Resultados generales de las aplicaciones piloto.}
  \label{cuadro:resultados}
  \footnotesize
  % Las columnas numéricas son p{} alineadas a la derecha, no r: así el
  % encabezado ("Inferencias registradas") puede partirse en dos líneas.
  % TeX no corta con guion la primera palabra de un párrafo, así que cada
  % columna tiene que caber su palabra más larga ("Requerimientos").
  \begin{tabularx}{\linewidth}{@{}%
      >{\raggedright\arraybackslash}X%
      >{\raggedright\arraybackslash}p{0.23\linewidth}%
      >{\raggedleft\arraybackslash}p{0.10\linewidth}%
      >{\raggedleft\arraybackslash}p{0.12\linewidth}%
      >{\raggedleft\arraybackslash}p{0.15\linewidth}%
      >{\raggedright\arraybackslash}p{0.12\linewidth}@{}}
    \toprule
    Piloto & Configuración & Viviendas & Inferencias registradas & Requerimientos cubiertos & Veredicto del sistema evaluado \\
    \midrule
    CASAS Aruba    & \texttt{config.yaml}         & 1 & 11\,175 & 9 de 9 & No cumple \\
    CASAS serie hh & \texttt{config.hogares.yaml} & 9 & 7561    & 9 de 9 & No cumple \\
    \bottomrule
  \end{tabularx}
  \fuente{elaboración propia con base en la ejecución documentada del repositorio.}
\end{table}

Los dos expedientes generaron la totalidad de los artefactos previstos: ficha de
caracterización, datasheet, protocolo de evaluación, explicaciones locales,
importancia global, resultados de desempeño desagregado, métricas de equidad,
bitácora de inferencias, reporte de explicabilidad, reporte de equidad, model
card, reporte de cumplimiento, bitácora de ejecución y manifiesto de hashes. En
ambos pilotos, los nueve requerimientos de la matriz quedaron respaldados por
evidencia presente y verificable.
% El manifiesto sella 14 artefactos: los de esta lista menos el propio
% manifiesto, más la figura de importancia global.

El resultado de no cumplimiento no indica un fallo del marco, sino un hallazgo
sobre el sistema evaluado. El procedimiento completo se ejecutó y produjo
evidencia suficiente para sostener el veredicto. Esto es consistente con el
criterio de éxito definido para el piloto: el marco debía demostrar que podía
generar un expediente de evidencia, incluso cuando la evidencia revelara
incumplimientos del sistema examinado.

En el módulo de explicabilidad, cada inferencia registrada en la bitácora quedó
vinculada con una explicación local. Esta relación satisface simultáneamente
R3.1 y R5.3, porque permite partir de una decisión registrada y llegar al
artefacto que explica las variables que contribuyeron a la predicción. La
agregación de esas explicaciones permitió construir una importancia global de
variables, asociada con R3.2. El uso de nombres de zonas del hogar permitió
generar enunciados interpretables para cuidadores y otros destinatarios no
técnicos, en concordancia con R3.3.

% Fuente: artefactos/reporte_explicabilidad.md y
% artefactos/hogares/reporte_explicabilidad.md (media de |SHAP|).
En Aruba, las variables con mayor importancia global, medida como la media del
valor absoluto de SHAP, fueron el movimiento en el sillón (0,0887), en la cocina
(0,0780) y en el dormitorio (0,0447), seguidas de la duración de la ventana y la
hora del día. En la configuración multi-hogar encabezaron la cocina (0,0452), el
baño (0,0440) y el dormitorio (0,0406). En ambos casos, las variables
dominantes son zonas del hogar, lo que permite contrastar el comportamiento del
modelo con el conocimiento del dominio sobre dónde ocurre cada actividad.

En el módulo de equidad, los resultados fueron desagregados por los subgrupos
declarados en la configuración. La existencia de umbrales sellados antes de
ejecutar las pruebas permitió diferenciar la medición de disparidades de una
interpretación posterior ajustada a los resultados. Cuando una combinación de
actividad, subgrupo y métrica no tuvo soporte suficiente, el marco la marcó como
no evaluable en lugar de contarla como aprobada. Esta decisión evitó convertir
ausencia de evidencia en evidencia de cumplimiento.

% Todas las cifras de esta parte salen de los expedientes: model_card.md,
% reporte_equidad.md y equidad/desempeno_desagregado.csv de cada piloto.
Como referencia para interpretar las disparidades, el modelo alcanzó una
exactitud de 0,704 y un F1 macro de 0,432 en Aruba, y de 0,326 y 0,276 en la
configuración multi-hogar, con una exactitud por vivienda entre 0,225 y 0,433.
Estas cifras no forman parte del criterio de éxito del estudio, pero condicionan
la lectura de los resultados de equidad: un clasificador de desempeño modesto
deja más margen para que aparezcan diferencias entre subgrupos.

En Aruba se evaluaron 32 de las 44 combinaciones de subgrupo, actividad y
métrica con efecto en el veredicto; las 12 restantes no alcanzaron el soporte
mínimo de 30 casos por tasa. Cuatro combinaciones superaron el umbral
(Cuadro~\ref{cuadro:aruba}). Las cuatro corresponden a la franja horaria y a la
diferencia de TVP: el modelo reconoció peor ciertas actividades según la hora
del día, mientras que ninguna combinación de tasa de falsos positivos ni de tipo
de día superó el umbral. Parte de esta disparidad puede deberse al diseño, pues
la franja se deriva de la hora, que es también una variable de entrada del
modelo. La mayor diferencia, en Sleeping, compara solo la madrugada y la noche,
porque la mañana y la tarde quedaron por debajo del soporte mínimo.

\begin{table}[htbp]
  \caption{Combinaciones que superan el umbral de 0,10 en Aruba.}
  \label{cuadro:aruba}
  \small
  \begin{tabularx}{\linewidth}{@{}>{\raggedright\arraybackslash}X*{5}{r}@{}}
    \toprule
    & & \multicolumn{4}{c}{TVP por franja horaria} \\
    \cmidrule(l){3-6}
    Actividad & Diferencia de TVP & Madrugada & Mañana & Tarde & Noche \\
    \midrule
    Sleeping         & 0,351 & 0,970 & 0,944* & 0,500* & 0,619 \\
    Otro             & 0,170 & 0,522 & 0,692  & 0,530  & 0,622 \\
    Meal\_Preparation & 0,169 & 0,880 & 0,870  & 0,798  & 0,712 \\
    Relax            & 0,128 & 0,835 & 0,901  & 0,773  & 0,820 \\
    \bottomrule
  \end{tabularx}
  \par\smallskip{\footnotesize * Menos de 30 casos: fuera de la comparación.\par}
  \fuente{reporte de equidad del expediente de Aruba.}
\end{table}

Entre las combinaciones no evaluables está Bed\_to\_Toilet, los traslados
nocturnos al baño: la prueba contiene 8 casos en la madrugada y ninguno en las
demás franjas. El marco no puede afirmar que el sistema reconozca por igual
esta actividad; sí verificó que sus falsas alarmas se mantuvieran dentro del
umbral en todas las franjas.

En la configuración multi-hogar se evaluaron 108 de 150 combinaciones y 29
superaron el umbral: 14 entre viviendas, 12 entre franjas horarias y 3 entre
tipos de día. A diferencia de Aruba, aquí los subgrupos corresponden a personas
distintas y no a condiciones de contexto, como plantea el requerimiento R4.1.
Las mayores diferencias entre viviendas se dieron en la TVP de Work (0,872),
Eat (0,810), Cook (0,503), Wash\_Dishes (0,447) y Sleep (0,425).

El caso de Cook muestra que la disparidad no se explica solo por la cantidad de
datos. La vivienda hh107, la de más casos de esa actividad en la prueba (243),
obtuvo la TVP más baja (0,292), mientras que hh103, con 122, obtuvo la más alta
(0,795). En Sleep, en cambio, hh107 superó a hh104 (0,867 frente a 0,442), las
dos únicas viviendas con soporte suficiente para esa actividad. El desempeño
desigual depende, así, de la combinación entre vivienda y actividad.
% D52 atribuye el caso de hh107 a que es la única vivienda con dos residentes.
% El dato no consta en el expediente ni en el crudo del repositorio.
\revisar{D52 atribuye la diferencia de hh107 a que tiene dos residentes; hace falta la fuente de ese dato para incluirlo}

En el módulo de trazabilidad, la bitácora estructurada permitió verificar
unicidad de identificadores, marcas temporales UTC, rango válido de confianza,
versión de modelo, versión de esquema, responsable declarado y existencia de la
explicación referenciada. La verificación final agregó la cobertura entre
conjunto de prueba y registros producidos, de modo que una bitácora válida pero
incompleta no pudiera sostener el cumplimiento.

Adicionalmente, el repositorio documentó 272 pruebas automatizadas en verde.
Estas pruebas cubrieron configuración, preparación de datos, modelado,
explicabilidad, equidad, trazabilidad, procedimiento y documentación. Aunque las
pruebas no sustituyen una auditoría externa, aportan evidencia de
reproducibilidad del procedimiento y reducen el riesgo de divergencia entre el
diseño normativo y la implementación computacional.
